% source: manti.01.ma, 1ère SM, série suites (Drive 0B6pfSJcEVng0WTQ5Uzh5MFlBbFE). Ex. 6 = RzBR Ex. 5 (non repris)
\ex{Suites}{WTQ5}{Ex. 1}{9}
Soit $(U_n)_{n\in\mathbb{N}}$ la suite définie par $U_0=1$ et $U_{n+1}=\dfrac{U_n}{3^nU_n+3}$. On pose, pour tout $n\in\mathbb{N}$, $V_n=\dfrac{1}{3^nU_n}$.
\begin{enumerate}
\item Calculer $U_1$, $V_0$ et $V_1$.
\item Montrer que $(V_n)_{n\in\mathbb{N}}$ est une suite arithmétique, puis calculer $V_n$ en fonction de $n$.
\item En déduire que $U_n=\dfrac{1}{3^{n-1}(n+3)}$.
\end{enumerate}
\fin
\ex{Suites}{WTQ5}{Ex. 2}{9}
Soit $(U_n)_{n\in\mathbb{N}}$ la suite définie par $U_0=1$ et $3U_{n+1}=2U_n+n+3$, et soit $(V_n)_{n\in\mathbb{N}}$ la suite définie par $V_n=U_n-n$.
\begin{enumerate}
\item Montrer que $(V_n)_{n\in\mathbb{N}}$ est une suite géométrique en précisant sa raison et son premier terme.
\item En déduire que $U_n=\left(\dfrac23\right)^n+n$.
\item Montrer que $U_0+U_1+\dots+U_{99}=4953-3\left(\dfrac23\right)^{100}$.
\end{enumerate}
\fin
\ex{Suites}{WTQ5}{Ex. 3}{10}
Soit $(U_n)_{n\in\mathbb{N}}$ la suite définie par $U_0=2$ et $U_{n+1}=\dfrac15\left(U_n-4n-1\right)$. On pose $V_n=U_n+n-1$ pour tout $n\in\mathbb{N}$.
\begin{enumerate}
\item \begin{enumerate}
\item Montrer que $(V_n)_n$ est une suite géométrique de raison $q=\dfrac15$.
\item Calculer $V_n$ et $U_n$ en fonction de $n$.
\end{enumerate}
\item On pose $T_n=V_0+V_1+\dots+V_n$ et $S_n=U_0+U_1+\dots+U_n$.
Montrer que $T_n=\dfrac14\left(5-\dfrac{1}{5^n}\right)$ et que $S_n=T_n-\dfrac{(n-2)(n+1)}{2}$.
\end{enumerate}
\fin
\ex{Suites}{WTQ5}{Ex. 4}{11}
Soit $(U_n)_{n\in\mathbb{N}}$ la suite définie par $U_0=0$, $U_1=1$ et $25U_{n+2}=10U_{n+1}-U_n$.
On pose, pour $n\in\mathbb{N}$, $W_n=5^nU_n$ et $V_n=U_{n+1}-\dfrac15U_n$.
\begin{enumerate}
\item Montrer que $(V_n)_n$ est une suite géométrique et que $(W_n)_n$ est une suite arithmétique.
\item Déterminer le terme général $U_n$ en fonction de $n$.
\item \begin{enumerate}
\item Montrer que $(\forall n\in\mathbb{N}^*)\ 0\leq U_{n+1}\leq\dfrac25U_n$.
\item En déduire que $(\forall n\in\mathbb{N}^*)\ 0\leq U_n\leq\left(\dfrac25\right)^{n-1}$.
\end{enumerate}
\end{enumerate}
\fin
\ex{Suites}{WTQ5}{Ex. 5}{11}
Soit $(U_n)_{n\in\mathbb{N}}$ la suite définie par $U_0=\dfrac{31}{5}$ et $U_{n+1}=\dfrac{U_n^2+5U_n+10}{3(U_n-1)}$.
\begin{enumerate}
\item Montrer que $(\forall n\in\mathbb{N})\ U_n\geq5$.
\item Montrer que $(U_n)_{n\in\mathbb{N}}$ est décroissante.
\item Montrer que pour tout $n\in\mathbb{N}$ : $U_{n+1}-5\leq\dfrac{1}{10}(U_n-5)$, et en déduire que $U_n-5\leq\dfrac65\left(\dfrac{1}{10}\right)^n$.
\end{enumerate}
\fin
\ex{Suites}{WTQ5}{Ex. 7}{10}
Soit $(U_n)_{n\in\mathbb{N}}$ la suite définie par $U_0=1$, $U_1=5$ et $U_{n+2}=4U_{n+1}-4U_n$. On pose $x_n=U_{n+1}-kU_n$.
\begin{enumerate}
\item Déterminer $k$ pour que $(x_n)_n$ soit une suite géométrique et préciser sa raison.
\item On suppose $k=2$.
\begin{enumerate}
\item Montrer que $U_{n+1}=2U_n+3\times2^n$.
\item On pose $U_n=2^nW_n$. Montrer que $(W_n)_n$ est une suite arithmétique.
\item Déterminer $U_n$ en fonction de $n$.
\end{enumerate}
\end{enumerate}
\fin
\ex{Suites}{WTQ5}{Ex. 8}{6}
Soit $(U_n)_{n\in\mathbb{N}}$ une suite telle que, pour tout $n\in\mathbb{N}$ : $S_n=U_0+U_1+\dots+U_{n-1}=\dfrac{n^2+n}{3}$ \ (l'original écrit $U_0+\dots+U_n$ [?]).
Calculer $U_0$ et $U_1$, montrer que $(U_n)_{n\in\mathbb{N}}$ est une suite arithmétique et déterminer $U_n$ en fonction de $n$.
\fin
